\section{The Notion of an Elite}

Due to the nature of blog posts, you may feel like someone watching a movie starting from the middle. Nevertheless, there is an architectonic plan as summarized:

\begin{itemize}


\item \textbf{Thesis}: Tradition is the normal state for human spiritual well-being.

\item \textbf{Hypothesis}: The Western world is in crisis due to its forgetting of Tradition.

\item \textbf{Antithesis}: The proper response, then, is to revolt against the modern world.

\item \textbf{Synthesis}: Some few men will undertake the effort to restore Tradition.
\end{itemize}


Rene Guenon's \emph{The Crisis of the Modern World} is the fundamental text that proves the hypothesis. Julius Evola's \emph{Revolt Against the Modern World} is the text of the antithesis. By understanding the thesis, hypothesis and antithesis, work can begin on the synthesis.

What follows can only be an outline, since all of the topics have been developed in full in other posts. The quotes are from Guenon's book.

Since we are concerned primarily about the West, that is where the focus will be. Moreover, we agree with Evola that Guenon is the master of the 20\textsuperscript{th} century. If you have any objections to Gornahoor, most likely they begin with an objection to Guenon. If you disagree with the thesis, there is nothing much here for you; there is no desire to convert anyone. You are part of the anti-Tradition, which most people today consider sane and normal.

The more insidious movements are the counter-tradition. They claim to be part of Tradition, or ``influenced'' by Tradition as it is usually put. People gravitate to it since they are still under the sway of the modern world, while pretending to be in revolt. Or else, the ``revolt'' is strictly on the material plane. There is no ``influence''; you have to decide to go ``all in'' or not at all.

\paragraph{The Western Tradition}


Both Guenon and Evola postulate a polar or Hyperborean origin of our Tradition. If you prefer, you can presume this to refer to a ``spiritual north''. It began with a Primordial Solar Monotheism\footnote{\url{https://gornahoor.net/?p=3205}}. From there, the next step was the religion that developed around the Vedas in India, then the paganisms in Europe, and finally the Catholicism of the Middle Ages. As St Augustine wrote, ``there is but one Tradition, which is now called Catholicism.'' As Guenon wrote:

\begin{quotex}
It is moreover quite certain that it is in Catholicism alone that all that may still remain of the traditional spirit in the West has been preserved.
\end{quotex}


Hence, that is where we start to begin the synthesis. Guenon confirms that:

\begin{quotex}
The least fantastic venture, in fact the only one that does not come up against immediate impossibilities, would therefore be an attempt to restore something comparable to what existed in the Middle Ages, with the differences demanded by modifications in the circumstances
\end{quotex}


The counter-tradition rejects that in favor of some imaginary past or future paganism. This is hardly necessary, since all that is valuable in paganism has been preserved, as we have repeatedly demonstrated. These are some highlights:

\begin{itemize}


\item The Church Fathers regarded Christianity as the esoteric aspect of Greek philosophy

\item Evola, no friend of Christianity, conceded that the Middle Ages preserved essential pagan elements.

\item Pope Benedict XVI made the point that once the Logos was admitted into the Gospel, the whole of Greek philosophy came along with it.
\end{itemize}


We can rightly extend that idea. Once Greek philosophy is becomes a part of Western Tradition, then we are justified to also pull in the Vedic Tradition. We have been justifying that point. Moreover, Guenon claims it is a necessary step:

\begin{quotex}
If the Catholic Church could thus be brought into touch with representatives of the Eastern traditions, it would be a preliminary step.
\end{quotex}


We have taken that step, and several others, relying on the likes of Bede Griffiths and Guido De Giorgio. Anyone still interested in debating ``paganism vs Christianity'' is wasting your time and has nothing to say in regard to the proposed synthesis.

\paragraph{Types of Religion}


Any discussion of the religions of the West needs to take into account three degrees of its manifestation, which correspond to Evola's three stages:

\begin{itemize}


\item Folk religion

\item Exoteric religion

\item Esoteric religion
\end{itemize}


The folk stage is tied up with superstition, the desire for material protection, unusual phenomena, and so on. The exoteric stage is roughly what is called ``organized religion''. Beyond doctrine, there is a lot of focus on ``churchy'' news and so on. We are not concerned with that here, while not rejecting it. A fortiori, we don't identify a Tradition solely with those in whom it is entrusted. What some bishop does or does not do or say is not relevant to the synthesis.

We don't need to dwell on the esoteric aspect here, since it has been covered many times. So, if someone wants to reject the Western tradition and replace it with some ``neo-paganism'', the obvious question, is which form? They deny the pagan folk religion with its belief in bodily gods and so on. They can try to convince a population to adopt some new exoteric form, but for what end? Ultimately, they need to deal with the esoteric aspect and, in that, no such replacement is possible. Perhaps a development or a deepening, but that can be accomplished more easily with the forms that currently exist.

The entire concept is pointless. As a personal example, there are two Indian colleagues with whom I have frequent conversations. When we discuss meditation, witnessing thoughts, staying watchful, and so on, there is mutual understanding, despite our different exoteric commitments.

\paragraph{Race and Tradition}


We can say a word about this, based on the interest in recent posts. First of all, the task is the return to Tradition. It is not to ``save'' some ill-defined ``Western civilization'', nor especially is the goal to ``save the white race'', assuming that is even necessary. Since Evola wrote more on this than did Guenon, we can fill in some details as he saw them.

Based on Bachofen's studies, Evola accepted the early existence of gynocratic and matriarchal societies in remote history or pre-history. However, the Hyperborean peoples never went through that stage. He attributed that to the masculine spirit of the Hyperboreans and their branches. However, even within the Hyperboreans, the feminine race has more in common with the gynocratic societies. The Hyperborean branches decline when they start adopting the feminine perspective. You can read his descriptions and decide if they match what you see.

One of the hallmarks is the lack of a sense of differentiation. For example, those who call for a united ``white race'' lack that sense. It simply is not a ``possibility of manifestation'' in Guenon's phrase. Were it a real possibility, it would have happened by now. The desire for unity on the material plane is a feminine trait, suitable for women or those beings with a man's body but a feminine spirit.

Men understand differentiation, since there are many kinds of people and ways of being in the world. There are various castes, for example, or ``races of the spirit'' as Evola calls them, and so on. There may be little in common between them, despite outward racial characteristics. For our purposes, especially, this brings up Guenon's ideal of the ``elite'', which is the greatest differentiation.

\paragraph{Intellectuality and the Elite}


Guenon is emphatic that Tradition can only be restored by an ``intellectual elite''. By necessity, an intellectual understanding is possible only for a few. Guenon writes:

\begin{quotex}
This is why a true understanding can come only from above and not from below; and this should be taken in a twofold sense: the work must begin from what is highest, that is, from principles, and descend gradually to the various orders of application, always keeping rigorously to the hierarchical dependence that exists between them; and it must also of necessity be the work of an elite in the truest and most complete meaning of this word: by this we mean exclusively an intellectual elite, and in reality, there can be no other.
\end{quotex}


Hence, the synthesis must begin with an understanding of principles, not from political concerns and most definitely not from biological concerns. Unfortunately, I cannot just tell you what such an elite would think. There are certainly texts and the best place to start is with your own Tradition. We have convincingly demonstrated that the Western tradition is still alive and viable.

As Guenon points out, such an elite is beyond all forms, so it doesn't help to look for another organization or formal order. Rather they will arise as needed. Meanwhile, there are certainly some things to do to prepare or to be able to recognize an elite. The first requirement is an intellectual conversion which represents the return to the center. For this, a sane and normal way of life, in its traditional understanding, is a prerequisite. Then, the practice of one's exoteric religion, for reasons that we've discussed many times.

If you choose to ignore opportunities in your life that come up, ask yourself where your life is headed. I watched a few episodes of a TV show called Dig, based on its premise. I once thought it may be worth reviewing, but that won't happen. However, there was one incident that raised this thought.

One of the characters is a boy who was being groomed to be the next High Priest of Jerusalem. One of his handlers, a prole woman and former junky, got it into her mind to rescue the boy from that, so she kidnapped him. It got me to wonder. On the one hand, he is on a high spiritual path (speaking figuratively, not that I am endorsing the specifics of the show). But what is she offering him against that, except a life of TV, video games, pop tarts, and so on?

One thing for certain is that the end of this world is a possibility of manifestation. Where will you be?


\flrightit{Posted on 2015-04-30 by Cologero}

\begin{center}* * *\end{center}

\begin{footnotesize}
\begin{sffamily}
\texttt{Name on 2015-05-01 at 03:51 said:}

Another great post, Mr. Salvo.

Paragraphs which I like:

``They claim to be part of Tradition, or ``influenced'' by Tradition as it is usually put. People gravitate to it since they are still under the sway of the modern world, while pretending to be in revolt. Or else, the ``revolt'' is strictly on the material plane. There is no ``influence''; you have to decide to go ``all in'' or not at all.

The folk stage is tied up with superstition, the desire for material protection, unusual phenomena, and so on. The exoteric stage is roughly what is called ``organized religion''.

A fortiori, we don't identify a Tradition solely with those in whom it is entrusted. What some bishop does or does not do or say is not relevant to the synthesis.

The desire for unity on the material plane is a feminine trait, suitable for women or those beings with a man's body but a feminine spirit.

For our purposes, especially, this brings up Guenon's ideal of the ``elite'', which is the greatest differentiation.

Hence, the synthesis must begin with an understanding of principles, not from political concerns and most definitely not from biological concerns.

The first requirement is an intellectual conversion which represents the return to the center.''

\hfill

\texttt{IA on 2015-05-01 at 10:31 said:}

I like a lot of what you say here but do have some problems with the term ``intellectual.'' Maybe a definition is in order. Or, if you've written about this or Evola or Guenon have you could refer me to such definition?

The problem is very well examined in Paul Johnson's ``The Intellectuals''.

\hfill

\texttt{Wayne Ferguson on 2015-05-01 at 11:35 said:}

Thesis: Tradition is the normal state for human spiritual well-being.

Hypothesis: The Western world is in crisis due to its forgetting of Tradition.

Antithesis: The proper response, then, is to revolt against the modern world

Synthesis: Some few men will undertake the effort to restore Tradition

That seems a rather odd use of the word ``antithesis''. It seems more like a minor premise or-- if we construe the hypothesis to be a premise --the conclusion to an augument in standard syllogistic form:

Major Premise: Tradition is the normal state for human spiritual well-being.

Minor Premise: The Western world is in crisis due to its forgetting of Tradition.

Conclusion: The proper response, then, is to revolt against the modern world

Corollary: Some few men will undertake the effort to restore Tradition

\hfill

\texttt{Wayne Ferguson on 2015-05-01 at 11:41 said:} 

``Both Guenon and Evola postulate a polar or Hyperborean origin of our Tradition. If you prefer, you can presume this to refer to a ``spiritual north''. It began with a Primordial Solar Monotheism. From there, the next step was the religion that developed around the Vedas in India, then the paganisms in Europe, and finally the Catholicism of the Middle Ages. As St Augustine wrote, ``there is but one Tradition, which is now called Catholicism.'' Why no mention of Guenon's tradition of choice? And what of Judaism as a tradition in its own right?

\hfill

\texttt{Cologero on 2015-05-01 at 12:43 said:}

Mr Ferguson, you must have gotten out of the wrong side of the bed this morning; that's the only explanation I can think of.

Perhaps I've taken some artistic liberties, but your objection is off the mark.

The hypothesis is the statement that needs to be proved. The proof, I contend, is in Guenon's Crisis of the West.

Re antithesis, this is from wikipedia:

\begin{quote}\footnotesize

According to Aristotle, the use of an antithesis makes the audience better understand the point one is trying to make through their argument. Further explained, the comparison of two situations or ideas makes choosing the correct one simpler. Aristotle states that antithesis in rhetoric is similar to syllogism due to the presentation of two conclusions within a statement.
\end{quote}


You are not the best audience, but the intended point is to emphasize the point. If the modern world is truly in crisis, then revolt is the proper response as claimed.

Synthesis is the act of building up.

As for your other ``point'', if I write this sentence:

\begin{quotex}\footnotesize

Since we are concerned primarily about the West, that is where the focus will be.
\end{quotex}


you can take that as my actual intention. Islam and Judaism are not products of the West, that is why there is no mention.

As for why Guenon ended up in Islam, that is certainly a topic we've addressed. If it is truly a burning issue for you, you can do the homework to look it up.

\hfill

\texttt{Cologero on 2015-05-01 at 12:47 said:}

Mr IA, if I write a sentence like this:

\begin{quote}\footnotesize

Guenon is emphatic that Tradition can only be restored by an ``intellectual elite''.
\end{quote}


that means that I am quoting Guenon. Since I even wrote this:

\begin{quotex}\footnotesize

The quotes are from Guenon's book.
\end{quotex}


obviously referring to the Crisis.

Did you bother to consult that work? Do you know what ``intellectual'' means in metaphysics? Did you ever consider that Johnson's use of the word is incorrect?

\hfill

\texttt{Wayne Ferguson on 2015-05-01 at 13:41 said:}

My apologies--I did not mean to sound so grumpy. Your explanation is sufficient, but the references to the Vedas and India left me wondering how the other traditions would fit in. My exposure to the thesis/antithesis/synthesis triad has, in the past, been that commonly attributed to Hegel... 

\hfill

\texttt{Alistair Fraser on 2015-05-01 at 17:21 said:}

This is a quite fabulous clarification , in my perception

\hfill

\texttt{Mark Citadel on 2015-05-01 at 18:27 said:}

I'm going to have to link to this upcoming. It's lucky Guenon's works are easily accessible, and I mean that in both senses of the word. The man comes off as very elegant and open to the layman as opposed to Evola who can often appear as an enigma.

I am in the preliminary stages of writing a book at this time, a sort of overview of the Reactionary ideal. Would it be okay to quote small sections of your work, attributed of course? I think you elucidate the non-conflict between Traditional Christianity and what Guenon espouses so well.

\hfill

\texttt{IA on 2015-05-01 at 20:35 said:}

``I'm going to have to link to this upcoming.''

I wouldn't if I were you. Too much the prima donna. They don't age well.

\hfill

\texttt{White Rabbit on 2015-05-01 at 23:43 said:}

``But what is she offering him against that, except a life of TV, video games, pop tarts, and so on?''

That's why young men are ``dropping out'' and finding their passion in playing video games all day. Tendies and pop tarts

\hfill

\texttt{David on 2015-05-02 at 13:01 said:}

Greetings. I saw the term « Christianity » and « Catholicism » used (sometimes) interchangeably on Gornahoor. Sometimes specific topics of Roman Catholicism (RCC) and Eastern Orthodoxy were addressed. But here the focus seems to be directed specifically towards the RCC : is it because you deem the Orthodoxy as a product of the East ? Is it because I am misunderstood and you therefore use the word Catholic as the ''One and Holy Catholic and Apostolic Church'' therefore considering the Orthodoxy ? Because on many topics the Orthodoxy is more akin to the Tradition that is possess, and since the Metropolitan and other Patriarch do not have (and didn't had as far as I know) the same power as the Pope, they do not possess the same Guelph pretension that the RCC has. Also let's include the fact that some of the Eastern European countries such as Greece, Romania, etc. are all Orthodox; and that Orthodoxy is claiming people from RCC and Protestantism (slowly of course) in European countries and America. Therefore I am asking you what to make of Eastern Orthodoxy, or if it is already within your definition of Catholic. Thank you.

\hfill

\texttt{Cologero on 2015-05-03 at 09:30 said:}

David, I don't see how the focus is directed exclusively toward the RCC, but rather the focus is on Tradition. As I've made clear, our true focus is on the \textit{Church of John}\footnote{\url{http://www.meditationsonthetarot.com/the-church-of-john}}, which is beyond form, yet not opposed to the exoteric church.

Apparently you are unaware that there is no more schism. The mutual ex-communications were lifted 20 years ago in 1995, even if there is not yet actually full ``communion''. With all the emphasis we've placed on the Philokalia, Cavarnos, Berdyaev, Solovyov, and Mouravieff, your objection makes no sense.

We are not interested in promoting or defending specific institutions ... we simply assume that one should participate in an exoteric tradition.

Guenon, and even Evola, use the Western Middle Ages as an example of Tradition, so that is as ``akin'' to Tradition as you could want. Certainly the West is more pagan than the East, because of the Frankish elements which tend more to legalism and objectivism. Nevertheless, the Germanics adopted the Roman religion.

The East never passed through a ``Middle Age'' ... it went from the Roman Empire to being overrun by Islam. So the differences are contingent, the result of historical events, but not essential, even if some prefer to battle indefinitely over the unessential.

The Eastern churches are far from an ideal model. They have their own problems with modernity ... for example, they boast about Tom Hanks despite his public scandal. There are many like him. In personal dealings with laypeople, I haven't seen any great love for or understanding of Tradition. In fact, they often confuse the specifically ethnic elements with transcendent truths.

So, if you think that switching allegiances, or the occasional poaching between Churches, is meaningful, we disagree. The fundamental issue still remains.

\hfill

\texttt{Jacob on 2015-05-03 at 16:48 said:}

I really like this post and can't figure out why anyone who understands the previous material would have any problems with it. The task has been laid out, but because it's been found difficult many don't want to even try it. It's a lot easier to clamor for regime change.

\hfill

\texttt{IA on 2015-05-03 at 17:24 said:}

Jacob,

If you are referring to me I will try to answer but as I mentioned this problem is dealt with by the devout, traditional Catholic artist and writer Paul Johnson.

Basically, the entire modernist world was created by intellectuals. Starting with Rousseau.

Roger Scruton also writes about intelligent but socially maladroit individuals who harbor resentments against powerful men.

I think both authors are concerned about the motivations of ``intellectuals.''

\hfill

\texttt{David on 2015-05-03 at 19:38 said:}

Cologero. I didn't know of the uplifting of 1995. In what kind of document was it made ? Thank you.

As far as the rest of the information, it goes without saying that what you say is logical. Yet my (visibly poorly formulated question) was more linked with your third paragraph ''We are not interested etc... '', so that answers my question : you use the ''real'' term of Catholic as the One church, and not as the RCC. Sorry for any misreading. Thank you.

\hfill

\texttt{Cologero on 2015-05-03 at 20:13 said:}

David, here is the document:

\textit{COMMON DECLARATION SIGNED IN THE VATICAN BY POPE JOHN PAUL II AND PATRIARCH BARTHOLOMEW I}\footnote{\url{http://www.ewtn.com/library/PAPALDOC/BARTHDEC.HTM}}

\hfill

\texttt{Matt on 2015-05-04 at 02:11 said:}

IA,

When Guenon -- and Evola too -- speaks of an intellectual elite, he doesn't have in mind what Johnson and Scruton write about. The intellectual elite is not composed of the common types from academia, or intelligent individuals that have chained themselves to their resentments; rather, it is composed of persons who directly and completely grasp the essences of things, which are synonomous with the divine ideas.archetypes, or to use Guenon's terminology, the possibilities of manifestation. Furthermore, they have the will to help carry out the actualization of the possibilites and to close the gap/privation that can arise between the concret particular and its respected archetype. On that second point, Evola goes into further writing about it than Guenon did.

\hfill

\texttt{IA on 2015-05-04 at 09:47 said:}

Matt,

Thank you for taking the time to offer a polite response.

I don't have a problem with anyone who has a grasp of transcendent archetypes. In fact, I rather like the idea if these persons are poets and artists or devout and humble monks and saints who can help others achieve the same experience.

What I do have trouble with is the idea of being coerced by an elite who actually want to block me from my own development. They do this not out of a love of knowledge but by love of a system of knowledge, two very different things.

\hfill

\texttt{Cologero on 2015-05-05 at 08:18 said:}

By law, Mr Citadel, you have the right to fair use for your own work. My concern is to be lumped in as part of some amorphous movement of ``neo-reaction'', ``new right'', etc.

\hfill

\texttt{Mark Citadel on 2015-05-05 at 14:01 said:}

I am neither part of the ENR nor do I label myself a NeoReactionary (I abandoned the label over concerns about certain elements which seemed to want to cherry-pick from Tradition as and when it suited the movement a la Mussolini's mistake).

I describe myself plainly as a Reactionary, working off a synthesis of various right wing thinkers, including De Maistre, Evola, Sarda, and others. I would not associate you intentionally with any specific group. My admiration for Gornahoor is in the intellectual rigor with which these theories are analyzed. You're a valuable source of information.

\hfill

\texttt{Br. Giles Mary on 2015-05-06 at 14:17 said:}

This is a fascinating post, Cologero.

Thesis: Tradition is the normal state for human spiritual well-being.

\textgreater Granted: We don't live in ``normal times''.


Hypothesis: The Western world is in crisis due to its forgetting of Tradition.

\textgreater Yes... Of course Guenon would say that it is one aspect of the manifestation of universal possibility.

Antithesis: The proper response, then, is to revolt against the modern world.

\textgreater Right now, I think I'd say, we should revolt against the anti-/counter-tradition.

I agreed with the basic premise of Evola's \textit{Ride the Tiger}, that the new disorders offer new possibilities. That means some kind of dialogue with the modern world and not a wholesale revolt. I don't think an elite would exist for itself, it should serve the brethren, who are by and large modern folk. Do you agree?

Synthesis: Some few men will undertake the effort to restore Tradition.

\textgreater This is interesting. The way I see it, the metaphysical principles are the most important instrument for countering the counter-tradition. We should act like men, not sitting by passively while error reigns. And a material unity is effeminate, like you said. And few are capable of thoroughgoing metaphysical knowledge. I think Guenon would say that a man with metaphysical aptitude need not resemble anything like what we'd call an intellectual. But, I think we should be concerned about what the bishops are doing. Many of them do a lot of good. Many do good in spite of themselves. And unfortunately, it seems, many have made nice with the modern world, falling for the hermeneutic of discontinuity, as Benedict XVI would say. I leave this quote from the Summa Theologiae. You might substitute, ``member of the faithful'' for ``religious'' if my point has any substance.

``I answer that, As Augustine says (Gen. ad lit. xii, 16), ``the agent is ever more excellent than the patient.'' Now in the genus of perfection according to Dionysius (Eccl. Hier. v, vi), bishops are in the position of ``perfecters,'' whereas religious are in the position of being ``perfected''; the former of which pertains to action, and the latter to passion. Whence it is evident that the state of perfection is more excellent in bishops than in religious.''

And from Lumen Gentium \#27: ``In virtue of this power, bishops have the sacred right and the duty before the Lord to make laws for their subjects, to pass judgment on them and to moderate everything pertaining to the ordering of worship and the apostolate.''

Certainly, Faith, Hope, and Love trump any status we might have within the Church... But, bishops have power to curb or assist anyone in their worship of God and service to the brethren... 

\hfill

\texttt{Wayne Ferguson on 2017-06-27 at 18:57 said:} 

``Pope Benedict XVI made the point that once the Logos was admitted into the Gospel, the whole of Greek philosophy came along with it.'' Might this be the source you have in mind?

``Every person as a rational being shares in the Logos, carrying within himself a ``seed'', and can perceive glimmers of the truth. Thus, the same Logos who revealed himself as a prophetic figure to the Hebrews of the ancient Law also manifested himself partially, in ``seeds of truth'', in Greek philosophy.

``Now, Justin concludes, since Christianity is the historical and personal manifestation of the Logos in his totality, it follows that ``whatever things were rightly said among all men are the property of us Christians'' (Second Apology of St Justin Martyr, 13: 4).

``In this way, although Justin disputed Greek philosophy and its contradictions, he decisively oriented any philosophical truth to the Logos, giving reasons for the unusual ``claim'' to truth and universality of the Christian religion. If the Old Testament leaned towards Christ, just as the symbol is a guide to the reality represented, then Greek philosophy also aspired to Christ and the Gospel, just as the part strives to be united with the whole.

``And he said that these two realities, the Old Testament and Greek philosophy, are like two paths that lead to Christ, to the Logos. This is why Greek philosophy cannot be opposed to Gospel truth, and Christians can draw from it confidently as from a good of their own. ''

\url{http://w2.vatican.va/content/benedict-xvi/en/audiences/2007/documents/hf_ben-xvi_aud_20070321.html}


\hfill

\end{sffamily}
\end{footnotesize}

